\section{Further questions}
\begin{enumerate}
\item Can the natural $K=R$ amplification and horizon-twelve proof be transferred to the signed cofinal chart without increasing the support of the literal power-oriented compiler? A correct transfer must re-establish every raw sign and carry bound before using its semantic power values.
\item Can the single anchor's comparison circuit be reduced while retaining a complete signed index proof? Eliminating its own certification would make the cap circular; finite agreement is not enough.
\item Can pooled index errors be partitioned into a small number of groups to improve the variable/support Pareto curve between the two delivered signed endpoints? Every candidate must be fully expanded and every group bounded before any claimed saving.
\item Can the positive one-square coprimality mechanism lower degree without losing its operation savings? The present transport majorants are deliberately explicit and prove positivity even when the decoded transport quotient is zero.
\item Can the nonzero-center two-root mechanism improve the signed carrier while keeping its auxiliary root-integrality and fixed-strong-pair interpolation order? In particular a new modulus must not be chosen after the pair it is supposed to constrain.
\item Can the ordinary calibration and each new arithmetic interface be assembled into one end-to-end kernel-checked endpoint? That is distinct unfinished formalization work, not a property implied by the supplied finite checks.
\end{enumerate}
The useful stopping criterion for another optimization round is a fully specified polynomial, an unbounded equivalence proof with explicit dependencies, an authentic coefficient tuple, exact metric evidence and a portable replay. Fewer symbols in an overview, a loose upper bound, or a favorable numerical sample is not such a result.
